\begin{abstract}
Drift calibration in head-mounted eye tracking is dominated by parametric
classical methods (similarity, polynomial, RBF, and combinations) fit per
session on a small set of fixation points. A recent line of work proposes
to refine the resulting predictions with a lightweight neural network that
ingests the predictions of multiple classical baselines and outputs an
adaptive correction. We replicate one such pipeline and find that the
reported gains arise entirely from a single line of code that subtracts the
regression target from the network's input feature, a leakage that
silently inflates the published numbers (e.g.\ from a corrected $30.5$~px
on JuDo1000 to the published $5.8$~px). After re-evaluating every variant
of the proposed design under an honest pipeline, we find that no neural
ensemble of classical calibrators reliably improves on the strongest
classical baseline alone. We then propose a simple alternative: after the
standard calibration phase, collect $K$ extra fixations as an
\emph{online recalibration context}, estimate the per-trial mean residual,
and shrink it by a coefficient $\lambda \in [0,1]^2$ predicted by a
$\sim$1k-parameter MLP from per-trial bias statistics (mean magnitude,
variance, $K$). Trained leave-one-subject-out on twelve participants and
evaluated on a thirteenth, our shrinkage estimator reduces mean L2 error
by $17\%$ at $K=12$ over the strongest classical baseline and $7\%$ over
the best globally-tuned fixed shrinkage. We release all code, data
splits, and per-experiment results.
\end{abstract}
